% ============================================================
% Selbständigkeitserklärung (zwingend, Handbuch S. 94)
% Wortlaut wörtlich aus dem Handbuch übernommen. Da die Probearbeit zu
% zweit verfasst wird, erklärt und unterschreibt jede Autorin / jeder
% Autor die Erklärung einzeln (die Vorlage verwendet "ich", nicht "wir").
%
% Muss handschriftlich unterschrieben werden -- nach dem Ausdrucken von
% Hand ausfüllen.
% ============================================================

\chapter*{Selbständigkeitserklärung}
\addcontentsline{toc}{chapter}{Selbständigkeitserklärung}

\noindent
Ich erkläre hiermit, dass ich
\begin{itemize}
  \item diese Arbeit ohne unerlaubte fremde Hilfe angefertigt habe.
  \item keine anderen als die angegebenen Quellen benutzt habe.
  \item sämtliche Stellen, die wörtlich aus fremden Quellen entnommen
    oder mit eigenen Worten wiedergegeben wurden, als solche
    gekennzeichnet und die Urheberschaft angegeben habe.
  \item die Meinungen und Ideen anderer explizit ausgewiesen habe.
  \item die Arbeit in eigenen Worten formuliert habe.
\end{itemize}

\noindent
Ich bin damit einverstanden, dass meine Arbeit als pdf-Datei in der
Mediothek archiviert wird und Schulangehörige diese einsehen dürfen.

% Bei Verwendung von KI-Tools (siehe ki-tools.tex) muss gemäss Handbuch
% Kap. 7.4.2 zusätzlich ein Reflexionsformular eingereicht werden --
% siehe reflexionsformular.tex (eigenständiges Dokument, jede/r Autor/in
% füllt ein eigenes Exemplar aus).

\vspace{3cm}

\noindent
\begin{minipage}{0.45\textwidth}
  Ort, Datum \\[2.5cm]
  \hrulefill \\
  Vorname Nachname 1
\end{minipage}
\hfill
\begin{minipage}{0.45\textwidth}
  Ort, Datum \\[2.5cm]
  \hrulefill \\
  Vorname Nachname 2
\end{minipage}
